%% Elements of Nested Fibrational Cosmology
%% Book VII: Canonical Closure Schema and Limits of the Spine
%% Final-Pass Canonical Draft
%%
%% Status legend:
%%   [U]  Unconditional  — proved from standing definitions and prior Unc results
%%   [C]  Conditional    — depends on explicitly declared hypotheses
%%   [D]  Definition-structural — structural, no proof obligation
%%   [R]  Remark only    — expository, non-load-bearing
%%   [O]  Open obligation — not yet proved

\documentclass[12pt,oneside]{amsbook}

\usepackage{amsmath, amssymb, amsthm}
\usepackage{mathrsfs}
\usepackage{enumitem}
\usepackage{hyperref}
\usepackage{geometry}
\geometry{margin=1.2in}

%% ---------------------------------------------------------------
%% Theorem environments
%% ---------------------------------------------------------------

\theoremstyle{plain}
\newtheorem{theorem}{Theorem}[section]
\newtheorem{lemma}[theorem]{Lemma}
\newtheorem{proposition}[theorem]{Proposition}
\newtheorem{corollary}[theorem]{Corollary}

\theoremstyle{definition}
\newtheorem{definition}[theorem]{Definition}
\newtheorem{standingrule}[theorem]{Standing Rule}

\theoremstyle{remark}
\newtheorem{remark}[theorem]{Remark}
\newtheorem{scholium}[theorem]{Scholium}

%% Status tag macro
\newcommand{\status}[1]{\textup{[#1]}}

\begin{document}

\title{Book VII\\[6pt]Canonical Closure Schema and Limits of the Spine}
\author{Elements of Nested Fibrational Cosmology}
\date{}
\maketitle

\tableofcontents

\bigskip

%% ---------------------------------------------------------------
\section*{VII.0\quad Purpose}
%% ---------------------------------------------------------------

Book~VII is the constitutional closure book of the canonical spine.
It does not add new source content, new branch endpoint content, or
new continuum interface content.  Its task is different and
indispensable: to regulate the force of all prior results, to govern
the lawful reuse of spine theorems in descendant books, to prohibit
silent promotion from scoped results to stronger claims, to preserve
named failure frontiers, and to distinguish architectural stability
from theorem-level completion.

Book~VII is therefore the force-governing book of the canon.  Without
it, a large source-descendant program of NFC's scope would drift into
ambiguity about what is proved, what is conditional, what is merely
suggestive, and what remains open.  The disciplines encoded here make
that drift impossible: every claim must declare its status, every
promotion must carry a transfer theorem, and every genuine obstruction
must remain visible as a named failure frontier.

Three prior texts supply the doctrinal content assembled here.  The
proof-status discipline of the spine --- each result stating its
hypotheses and listing only its licensed dependencies --- is the source
of Section~\ref{sec:status}.  The scoped-certificate and promotion law
--- a soundly scoped finite certificate carries force only in its
declared regime, and promotion beyond that regime requires a separate
transfer theorem --- is the source of Section~\ref{sec:promotion}.
The governance program's explicit treatment of named failure frontiers,
stable extensibility, and external falsifiability supplies
Sections~\ref{sec:frontiers} and \ref{sec:extensibility}.  Book~VII
makes all of this constitutional law.

%% ---------------------------------------------------------------
\section{Status and Dependency Law}
\label{sec:status}
%% ---------------------------------------------------------------

The first task of Book~VII is to make theorem force explicit and
immutable.

\begin{standingrule}[\status{D}]\label{sr:declare-status}
Every theorem, proposition, and corollary must declare:
\begin{enumerate}[label=\textup{(\arabic*)}]
  \item its status (unconditional, conditional, bridge, or open),
  \item its dependencies, and
  \item its scoped regime of force.
\end{enumerate}
\end{standingrule}

\begin{definition}[\status{D}]\label{def:unconditional}
An \emph{unconditional theorem} is a result proved from the standing
definitions, common notions, and previously certified unconditional
spine results alone, with no extra hypotheses.
\end{definition}

\begin{definition}[\status{D}]\label{def:conditional}
A \emph{conditional closure} is a result whose force depends on one or
more explicitly declared hypotheses not yet absorbed into the
unconditional spine.
\end{definition}

\begin{definition}[\status{D}]\label{def:bridge-theorem}
A \emph{bridge theorem} is a theorem carrying force from one scoped
regime to another: for example, from a certified finite regime to a
continuum target, or from source-side canonicality to a downstream
endpoint interpretation.
\end{definition}

\begin{definition}[\status{D}]\label{def:open-obligation}
An \emph{open obligation} is a named burden not yet discharged by
theorem or certified bridge.  Open obligations are mathematical
objects, not admissions of failure: a precisely stated open obligation
constrains future work and prevents false closure.
\end{definition}

\begin{proposition}[\status{U}]\label{prop:cite-declared-force}
A theorem may be cited only with the force licensed by its declared
status and dependencies.
\end{proposition}

\begin{proof}
A theorem's force is determined by the hypotheses and prior results on
which its proof depends.  If one cites the theorem with stronger force
than its declared dependencies support, one has asserted a stronger
claim without providing the stronger proof.  Therefore any citation
must respect the declared status.
\end{proof}

\begin{corollary}[\status{U}]\label{cor:no-unconditional-by-proximity}
No conditional result becomes unconditional by repetition, compression,
or proximity to stronger theorems.
\end{corollary}

\begin{proof}
The force of a result is determined by its proof, not by its context.
Repeating, compressing, or placing a conditional result next to a
stronger one does not discharge its extra hypotheses.  Therefore the
result remains conditional.
\end{proof}

\begin{remark}[\status{R}]\label{rmk:status-not-decoration}
Dependency discipline is not editorial decoration.  In a canon of this
scope, where theorems span seven books and dozens of derived branches,
the status tags [U], [C], [B], [O] are load-bearing structural
information.  A reader who misreads a [C] as a [U] has silently assumed
an unproved hypothesis.  Standing Rule~\ref{sr:declare-status} prevents
this by requiring explicit declaration.
\end{remark}

%% ---------------------------------------------------------------
\section{Scoped Certificates and Promotion Law}
\label{sec:promotion}
%% ---------------------------------------------------------------

\begin{definition}[\status{D}]\label{def:scoped-cert}
A \emph{scoped certificate} is a certified result whose force is
limited to the explicit regime named in its statement.
\end{definition}

\begin{definition}[\status{D}]\label{def:promotion}
A \emph{promotion} is any attempt to cite a result beyond the scoped
regime in which it was proved.
\end{definition}

\begin{definition}[\status{D}]\label{def:transfer-theorem}
A \emph{transfer theorem} for a promotion is a separately proved
theorem carrying a scoped result from its original regime into a
larger target regime.
\end{definition}

\begin{theorem}[\status{U}]\label{thm:scoped-cert-rule}
\textup{(Scoped Certificate Rule.)}  A scoped certificate proves only
the claim stated on its declared regime and carries no automatic
extension beyond that regime.
\end{theorem}

\begin{proof}
A scoped certificate is, by Definition~\ref{def:scoped-cert}, a proof
whose force is limited to its declared regime.  Any extension beyond
that regime is a different, stronger claim.  The proof of the scoped
certificate contains no argument for the stronger claim, so the
stronger claim is not established.  Therefore no automatic extension
follows.
\end{proof}

\begin{theorem}[\status{U}]\label{thm:promotion-requires-transfer}
\textup{(Promotion Requires Transfer Theorem.)}  A result may be
promoted beyond its scoped regime only when both of the following hold:
\begin{enumerate}[label=\textup{(\arabic*)}]
  \item a soundly scoped certificate establishes the result on the
        original declared regime, and
  \item a separately proved transfer theorem carries that certified
        content to the promoted target scope.
\end{enumerate}
Neither ingredient alone suffices.
\end{theorem}

\begin{proof}
A scoped certificate alone proves only the original claim on its
declared regime (Theorem~\ref{thm:scoped-cert-rule}).  A transfer
theorem alone cannot promote content that has not yet been certified
on its source regime: there is nothing to carry.  Therefore both
ingredients are necessary.  When both are present, the scoped
certificate provides the certified source content and the transfer
theorem carries it to the promoted target.  Together they constitute
a complete proof of the promoted claim.
\end{proof}

\begin{corollary}[\status{U}]\label{cor:finite-no-continuum}
Finite certification does not self-promote to continuum force.
\end{corollary}

\begin{proof}
A finite scoped certificate is limited to its declared finite regime
(Theorem~\ref{thm:scoped-cert-rule}).  Promotion to continuum force
requires a transfer theorem (Theorem~\ref{thm:promotion-requires-transfer}).
No such transfer is automatic.
\end{proof}

\begin{corollary}[\status{U}]\label{cor:source-no-endpoint}
Source-side canonicality does not self-promote to physical endpoint
identification.
\end{corollary}

\begin{corollary}[\status{U}]\label{cor:conditional-no-unconditional}
Branch-conditional success does not self-promote to unconditional
closure.
\end{corollary}

\begin{scholium}[\status{R}]\label{sch:promotion-law-universal}
Theorems~\ref{thm:scoped-cert-rule} and
\ref{thm:promotion-requires-transfer} govern every upward move in the
canon.  The promotion law is not limited to continuum promotion: it
governs promotion from finite to asymptotic, from source-side to
endpoint, from conditional to unconditional, and from one branch
regime to another.  This theorem family is among the deepest
constitutional laws of the program.
\end{scholium}

%% ---------------------------------------------------------------
\section{Status Inflation}
%% ---------------------------------------------------------------

\begin{definition}[\status{D}]\label{def:status-inflation}
\emph{Status inflation} is the illicit promotion of a claim beyond
the force licensed by its declared dependencies, scoped certificate,
and transfer stack.
\end{definition}

\begin{standingrule}[\status{D}]\label{sr:no-inflation}
No silent status inflation is permitted.
\end{standingrule}

\begin{standingrule}[\status{D}]\label{sr:precanonical-force}
\textup{(Pre-Canonical Remarks Have No Canonical Force.)}
Remarks in canonical spine books or branch books that reference
pre-canonical documents (retired monographs, dream-paper suites,
pre-canonical derivation records) have the status of archival
annotations only.  No result stated or sketched in a pre-canonical
document acquires canonical theorem-bearing force by being cited in
a canonical remark.  The canonical status of every claim is
determined solely by its proof or declared derivation within the
canonical papers.  A pre-canonical remark may be cited for:
\begin{enumerate}[label=\textup{(\roman*)}]
  \item identifying the route along which a future canonical proof
    should proceed,
  \item recording that a result is likely provable (providing
    heuristic confidence), or
  \item documenting the pre-canonical derivation as a historical
    record.
\end{enumerate}
A pre-canonical remark may \textbf{not} be cited as establishing,
proving, or discharging any open canonical obligation.  Any remark
of the form ``pre-canonical record contains a derivation of~$X$''
is equivalent to: ``$X$ is unproved in canonical papers and the
canonical obligation remains open; a candidate proof route exists.''
\end{standingrule}

\begin{standingrule}[\status{D}]\label{sr:proof-grade}
\textup{(Proof-Grade Classification.)}
Every theorem, lemma, or proof package in the canonical program should
be assigned a maturity level from the following scale, supplementing
the status tags $[\mathrm{U}]$/$[\mathrm{C}]$/$[\mathrm{O}]$/$[\mathrm{D}]$
with a finer internal-rigor assessment:
\begin{itemize}
  \item \textbf{P0}~--- Vision or heuristic only; no formal definitions,
    no proof.
  \item \textbf{P1}~--- Formal definitions exist; proof incomplete.
  \item \textbf{P2}~--- Proof scaffold exists; some lemmas structural
    rather than fully analytic.
  \item \textbf{P3}~--- Theorem package internally complete at
    CMP-like standard.
  \item \textbf{P4}~--- Externally paper-ready; referee-survivable.
\end{itemize}
Intermediate grades (e.g.\ P2.5, P2.7) may record partial progress.
A result may be $[\mathrm{C}]$ at~P3 (conditional but rigorously proved
under its hypotheses) or $[\mathrm{U}]$ at~P2 (claimed unconditional,
proof not yet fully assembled).  The proof-grade is recorded in the
Canon Ledger and governs priority for proof-hardening work.

Each result also carries a scope tag: \emph{source-level},
\emph{post-source/branch-neutral}, \emph{branch-specific},
\emph{continuum-interface}, or \emph{endpoint-specific}.  Scope tags
prevent contamination of branch-neutral results by branch-specific
assumptions.
\end{standingrule}

\begin{standingrule}[\status{D}]\label{sr:cmp-standard}
\textup{(Program-Wide Proof-Discipline Requirements.)}
Every layer of the canonical program is written to the following six
requirements:
\begin{enumerate}[label=\textup{(\arabic*)}]
  \item \textbf{Every object has a precise ambient home.}  Maps must
    have domain and codomain; norms, topologies, and equivalence
    relations must be stated.
  \item \textbf{Every theorem has a sharp status.}  Status tags are
    enforced together with exact hypotheses, exact dependencies,
    scoped regime, and a named failure frontier when the result is not
    unconditional.
  \item \textbf{Every selection principle must be formalized.}
    Phrases such as ``washed out under transport'' or ``absorbed under
    quotient'' are inadmissible unless converted into a definition,
    lemma, or theorem.
  \item \textbf{Branch-neutral results are separated from
    branch-specific interpretation.}  Every result carries a scope
    tag as defined in Standing Rule~\ref{sr:proof-grade}.
  \item \textbf{Heuristic ladders are split.}  Each development is
    explicitly partitioned into: proved core, formal conjectural
    extension, and roadmap material.
  \item \textbf{Every major theorem package declares its referee
    attack surface:} weakest assumptions, possible objections, where
    the proof burden lives, and what would falsify or downgrade the
    claim.
\end{enumerate}
Standing Rules~\ref{sr:declare-status}, \ref{sr:no-inflation}, and
\ref{sr:precanonical-force} are particular instances of requirements
(2) and (5).
\end{standingrule}

\begin{proposition}[\status{U}]\label{prop:downstream-force}
A downstream claim carries only the force of the named bridge stack
on which it depends.
\end{proposition}

\begin{proof}
A downstream claim's proof derives from a stack of prior results.  If
any result in that stack is conditional, scoped, or bridge-dependent,
the downstream claim inherits those limits: a conclusion is no stronger
than the weakest assumption it requires.  Therefore the downstream
claim carries only the force of its named bridge stack.
\end{proof}

\begin{corollary}[\status{U}]\label{cor:failed-bridge-reopens}
If a bridge theorem fails or remains open, every stronger downstream
claim depending on it is reopened.
\end{corollary}

\begin{proof}
By Proposition~\ref{prop:downstream-force}, the downstream claim
carries only the force of its bridge stack.  If one element of that
stack is removed (because the bridge fails) or unknown (because it is
open), the stack no longer supports the claimed conclusion.
\end{proof}

\begin{corollary}[\status{U}]\label{cor:slogan-no-force}
Compression of a proof chain into a slogan does not enlarge its
force.
\end{corollary}

\begin{proof}
The force of a claim is determined by its proof, not its verbal form.
Compressing a conditional proof chain into a confident-sounding
declaration does not discharge the conditions.
\end{proof}

\begin{remark}[\status{R}]\label{rmk:slogan-protection}
Corollary~\ref{cor:slogan-no-force} is one of the main protections
against premature closure rhetoric in a large source-descendant
program.  As the number of books and branches grows, there is
increasing pressure to summarize results in shorthand.  The
corollary insists that however a result is summarized, its underlying
status remains what the proof licenses.
\end{remark}

%% ---------------------------------------------------------------
\section{Named Failure Frontiers}
\label{sec:frontiers}
%% ---------------------------------------------------------------

A rigorous program needs an explicit theory of honest incompletion.

\begin{definition}[\status{D}]\label{def:named-frontier}
A \emph{named failure frontier} is a precisely formulated obstruction
identifying the exact missing ingredient that prevents a stronger
closure claim.  A named failure frontier is a positive mathematical
object: it records structural information about what is missing and
constrains future work.
\end{definition}

\begin{proposition}[\status{U}]\label{prop:frontier-valuable}
A descendant program may remain canonically valuable even when closure
fails, provided its failure frontier is explicit.
\end{proposition}

\begin{proof}
A named failure frontier records precisely what is missing.  This
information constrains future work (it says what must be supplied),
prevents false closure (it says what has not been proved), and
stabilizes the remaining burden (it says where the true difficulty
lies).  These are positive mathematical contributions, independent of
whether the closure is eventually achieved.
\end{proof}

\begin{theorem}[\status{U}]\label{thm:frontier-stability}
\textup{(Named Failure Frontier Theorem.)}  If a descendant branch is
blocked by a named failure frontier, then no equivalent restatement of
the same closure claim bypasses that obstruction without supplying the
missing structure.
\end{theorem}

\begin{proof}
A named failure frontier identifies a structural obstruction: a
specific ingredient that the closure claim requires but whose proof
is absent.  Any restatement of the same closure claim --- however
rephrased, reframed, or reterminologized --- still requires the same
structural content to be true.  Since the missing ingredient is
structural rather than verbal, rephrasing cannot supply it.  Therefore
no equivalent reformulation bypasses the frontier unless the missing
ingredient is actually provided.
\end{proof}

\begin{corollary}[\status{U}]\label{cor:rephrasing-not-progress}
Rephrasing is not progress unless the named obstruction has been
discharged.
\end{corollary}

\begin{definition}[\status{D}]\label{def:toolkit-boundary-frontier}
\textup{(Toolkit-Boundary Frontier.)}
A named failure frontier (Def.~\ref{def:named-frontier}) is a
\emph{toolkit-boundary frontier} when its missing ingredient is not
a gap in the branch's own construction but a class of question lying
outside the current canonical methods --- the collar-algebra,
transport-ledger, MSC/quotient-visibility, and continuum-bridge
machinery.  A toolkit-boundary frontier is distinguished from an
\emph{in-toolkit frontier}, whose missing ingredient is in
principle supplied by an instance of an existing canonical method,
and from a \emph{gated frontier}, whose discharge waits on a
specific named obligation already located in another branch.
\end{definition}

\begin{definition}[\status{D}]\label{def:frontier-residual-isolation}
\textup{(Residual Isolation.)}
A frontier is \emph{residual-isolated} when the branch has certified
every part of the original obligation that its methods reach, and
has reduced the remaining burden to an explicitly characterized
residual together with a proof that the residual is the unique
surviving case.  Residual isolation converts a broad frontier into a
sharply bounded one without flipping the obligation's status.
\end{definition}

\begin{proposition}[\status{U}]\label{prop:toolkit-boundary-no-inflation}
Certifying the in-toolkit portion of an obligation and isolating a
toolkit-boundary residual does not discharge the obligation: the
obligation retains [O] status until the residual itself is supplied.
\end{proposition}

\begin{proof}
By Def.~\ref{def:open-obligation} an obligation is discharged only
when its full closure claim is established at the declared force.
Residual isolation (Def.~\ref{def:frontier-residual-isolation})
establishes the closure claim on all cases except the isolated
residual; by Prop.~\ref{prop:frontier-valuable} the certified
portion and the sharpened frontier are positive contributions, but
the closure claim as a whole remains unproved while the residual
stands.  By the no-inflation standing rule
(Standing Rule~\ref{sr:no-inflation}), partial establishment may
not be recorded as full closure.  Hence the obligation retains [O].
The status tag therefore carries, for a residual-isolated
toolkit-boundary frontier, a different operational meaning than a
fresh undifferentiated frontier: not ``the branch has not yet
worked on this'' but ``the branch has reduced this to a
characterized residual outside current methods.''  The distinction
is recorded in the obligation text, not in the status tag, which
remains [O] in both cases.
\end{proof}

\begin{standingrule}[\status{D}]\label{sr:toolkit-boundary-honesty}
\textup{(Toolkit-Boundary Honesty.)}
When a branch reduces an open obligation to a toolkit-boundary
residual, it must (i) state the residual precisely, (ii) prove it is
the unique surviving case (residual isolation), and (iii) name the
class of method the residual would require, so that the frontier is
not mistaken for either an in-toolkit gap (which would invite a
spurious near-term discharge claim) or a vague incompletion.  A
toolkit-boundary frontier so recorded satisfies the named-failure
discipline of this section at the highest available precision.
\end{standingrule}

\begin{remark}[\status{R}]\label{rmk:toolkit-boundary-instances}
The current corpus carries four residual-isolated toolkit-boundary
frontiers, each reduced to a characterized residual by its branch
and retained at [O] under
Prop.~\ref{prop:toolkit-boundary-no-inflation}: the RH arithmetic
positivity frontier (\texttt{ob:rh-s1-formal}, residual: $Q[\psi]
\ge 0$ from the arithmetic structure of $\Xi$); the crystallographic
phase frontier (\texttt{ob:cryst-PHASE}, residual: phase uniqueness
modulo the certified invariance group for generic non-centrosymmetric
configurations); and the biological boundary-forcing frontier
(\texttt{ob:bio-BND-open}, residual: constructibility of a
coreless-complement persistent replicator).  Each names a residual
of structure-factor-fibration, arithmetic-positivity, or
configuration-existence character that the collar-algebra and
transport-ledger toolkit does not reach.  Distinct from these are
the \emph{gated} frontiers
(Def.~\ref{def:toolkit-boundary-frontier}): the SPEC matter-coupling
obligations, blocked not by method but by the named YM-side gate
O\_ID\textsuperscript{cont}.  This stratification --- toolkit-boundary
versus gated --- is itself governance information: it tells the
reader that no presently open obligation is both in-toolkit and
actionable by existing discharge methods, so further progress
requires either new methods (the toolkit-boundary residuals) or
discharge of the upstream gate (the gated obligations).
\end{remark}

\begin{remark}[\status{R}]\label{rmk:frontier-not-embarrassment}
A named failure frontier is not an embarrassment to be hidden or
minimized.  It is a mathematically valuable constitutional object.
Naming the frontier precisely is more informative than vague claims of
incompleteness: it tells the reader exactly what is needed, enables
targeted proof effort, and prevents the program from claiming more
than it has established.  The NS branch book, for example, is governed
by two explicit frontier objects (Ren.5 and SB2) rather than a generic
``NS regularity remains open.''  That precision is the discipline
Book~VII mandates.
\end{remark}

%% ---------------------------------------------------------------
\section{Stable Extensibility and External Falsifiability}
\label{sec:extensibility}
%% ---------------------------------------------------------------

A source-centered canonical program must support controlled growth
while remaining vulnerable to genuine failure.

\begin{definition}[\status{D}]\label{def:stably-extensible}
A descendant program is \emph{stably extensible} if new lawful
branches can be added by the same declared source-to-endpoint intake
law --- the branch legitimacy law of Book~V --- without requiring
structural changes to the source spine itself.
\end{definition}

\begin{definition}[\status{D}]\label{def:externally-falsifiable}
A descendant branch is \emph{externally falsifiable} if its claims
can, in principle, be rejected by failure of one of its declared
bridge steps, endpoint criteria, or source-to-target comparison
theorems.
\end{definition}

\begin{theorem}[\status{U}]\label{thm:stable-extensibility}
\textup{(Stable Extensibility Theorem.)}  The canonical spine is
stably extensible precisely when new descendant branches are admitted
only through the branch-legitimacy law of Book~V and the canonical
closure schema of this book.
\end{theorem}

\begin{proof}
If new descendants are admitted only through the fixed intake and
closure laws of Books~V and VII, then the source spine need not be
rebuilt for each new branch: $\mathbf{U}$, the POT category, and
Books~I--IV are unchanged.  Conversely, if branch admission requires
revising the source spine itself, then the spine is not stable under
the addition of new branches.  Therefore stable extensibility holds
exactly when branch admission is governed by the existing descendant
laws without structural modification to the source.
\end{proof}

\begin{theorem}[\status{U}]\label{thm:external-falsifiability}
\textup{(External Falsifiability Theorem.)}  A descendant branch is
externally falsifiable only if its spine imports, bridge stack,
endpoint target, and failure frontier are explicitly declared.
\end{theorem}

\begin{proof}
A claim is externally falsifiable only when the conditions under which
it would fail are identifiable by an independent reader.  If any of
the branch's spine imports, bridge stack, endpoint target, or failure
frontier are hidden or undeclared, an independent reader cannot
determine which specific step would need to fail to refute the claim.
Therefore explicit declaration of these items is necessary for
external falsifiability.
\end{proof}

\begin{corollary}[\status{U}]\label{cor:opaque-inadmissible}
Opaque branch claims --- those whose support structure is not fully
declared --- are not canonically admissible.
\end{corollary}

\begin{remark}[\status{R}]\label{rmk:falsifiability-honesty}
In a program of NFC's scope, falsifiability increases with
explicitness.  Hidden support always weakens the work, not because
hidden support is necessarily wrong, but because it cannot be examined,
challenged, or discharged.  A branch that declares all its support
structure openly is not only more honest; it is more scientifically
serious.
\end{remark}

%% ---------------------------------------------------------------
\section{Import Discipline and the Canonical Closure Schema}
%% ---------------------------------------------------------------

\begin{definition}[\status{D}]\label{def:admissible-import}
An \emph{admissible import} is a theorem from the canonical spine or
another properly certified source whose status, scope, and dependency
conditions are preserved under reuse.  An admissible import is cited
with its declared status intact; it is not promoted beyond that status
in the citing book.
\end{definition}

\begin{definition}[\status{D}]\label{def:prohibited-import}
A \emph{prohibited import} is any undeclared assumption, heuristic
gloss, or external structure used as though it had theorem-level force
without being listed in the branch's declared import ledger.
\end{definition}

\begin{definition}[\status{D}]\label{def:variable-licensing}
\textup{(Variable Licensing Predicate.)}
A \emph{physics variable} $X$ (energy, mass, charge, time,
temperature, distance, and the like) is \emph{licensed} within a
declared regime when both of the following hold:
\begin{enumerate}[label=\textup{(\alph*)}]
  \item\label{vl:form} $X$ is realized as at least one of:
  \begin{enumerate}[label=\textup{(\roman*)}]
    \item a quotient-visible invariant of the declared
          observational quotient;
    \item a transport-stable ledger in the sense of Book~III;
    \item a certified branch-interface coefficient with declared
          dependency conditions;
    \item a continuum encoding, via a separately proved Book~VI
          bridge, of certified discrete structure;
  \end{enumerate}
  \item\label{vl:decl} its scope, descent path, status tag, and
        anti-smuggling audit (the standard structures explicitly
        \emph{not} assumed) are declared.
\end{enumerate}
Use of a physics variable as a primitive --- outside
clauses~\ref{vl:form} and \ref{vl:decl} --- is a prohibited import
in the sense of Definition~\ref{def:prohibited-import}.  No physics
variable is fundamental by default.
\end{definition}

\begin{proposition}[\status{U}]\label{prop:variable-licensing-conservative}
\textup{(Conservativity of the Variable Licensing Predicate.)}
Definition~\ref{def:variable-licensing} introduces no new licensing
power and no new restriction: it is a consolidating restatement of
discipline already in force.  Specifically:
\begin{enumerate}[label=\textup{(\arabic*)}]
  \item every use of a variable satisfying the predicate was
        already admissible under the standing rules; and
  \item every use of a variable violating the predicate was
        already inadmissible.
\end{enumerate}
\end{proposition}

\begin{proof}
For (1), each disjunct of clause~\ref{vl:form} is individually
governed by an existing canonical instrument, and clause~\ref{vl:decl}
adds only declarations already required.  Quotient-visible invariants
are admissible exactly as observable content of the declared quotient;
transport-stable ledgers are governed by the Book~III transfer
machinery (the Observational Transfer Theorem of the Spine,
Book~III); branch-interface
coefficients with declared dependencies are admissible imports under
Definition~\ref{def:admissible-import} together with the Dependency
Declaration Theorem (Thm.~\ref{thm:dependency-declaration}); and
continuum encodings are licensed exactly by the Book~VI bridge schema
(Thm.~\ref{thm:continuum-bridge-schema}, Book~VI), which already
requires a separately proved bridge.  The declaration requirements of
clause~\ref{vl:decl} restate Standing
Rule~\ref{sr:declare-status} and
Thm.~\ref{thm:dependency-declaration}.  Hence the predicate licenses
nothing new.

For (2), suppose a variable is used in a manner violating the
predicate.  If it fails clause~\ref{vl:form}, then it enters as a
structure that is neither quotient-derived, transport-governed,
dependency-declared, nor bridge-encoded --- that is, as an undeclared
external structure carrying theorem-level force, which is a
prohibited import by Definition~\ref{def:prohibited-import}.  If it
fails clause~\ref{vl:decl}, the use violates Standing
Rule~\ref{sr:declare-status} or
Thm.~\ref{thm:dependency-declaration} directly.  Hence the predicate
forbids nothing that was previously permitted.

Both directions hold by case analysis over already-canonical
instruments; the predicate is therefore conservative.
\end{proof}

\begin{remark}[\status{R}]\label{rmk:variable-licensing-usage}
\textup{(Usage note.)}
Definition~\ref{def:variable-licensing} supplies the single citable
predicate underlying the variable-recovery shells maintained in the
speculative holding register (entropy, time, energy, Hamiltonian,
mass, charge, phase/frequency, metric/distance, curvature, momentum,
spin, temperature, couplings).  Spine-level instances already in
canon include the charge recovery shell
(Book~III, Rmk.~\ref{rmk:charge-recovery-shell}) and the SM matter
datum discipline (SM Branch).  Working-maturity vocabulary attached
to such shells in the holding register (\emph{theorem-shell},
\emph{conditional close}, \emph{template-backed},
\emph{near-closed}) is holding-internal and carries no canonical
force; no recovery shell may be cited at [U] or [C] except through
formal promotion (Standing Rule~\ref{sr:precanonical-force},
Thm.~\ref{thm:promotion-requires-transfer}).
\end{remark}

\begin{theorem}[\status{U}]\label{thm:dependency-declaration}
\textup{(Dependency Declaration Theorem.)}  Every descendant theorem
has canonical force only when its spine imports, extra hypotheses, and
bridge theorems are explicitly declared.
\end{theorem}

\begin{proof}
A descendant theorem's force depends on the precise support from which
it is derived.  If imports, added hypotheses, or bridge theorems are
undeclared, then the conditions under which the theorem holds cannot
be tracked by an independent reader, and the theorem cannot be cited,
challenged, or discharged.  Therefore explicit declaration is
necessary for canonical force.
\end{proof}

\begin{theorem}[\status{U}]\label{thm:closure-schema}
\textup{(Canonical Closure Schema Theorem.)}  Every derived branch
book must declare:
\begin{enumerate}[label=\textup{(\arabic*)}]
  \item its spine imports and the status of each imported result,
  \item its branch-specific inputs and hypotheses,
  \item its endpoint target,
  \item its bridge theorems,
  \item its named failure frontier, if any,
\end{enumerate}
and may claim no stronger status than these declarations jointly
permit.
\end{theorem}

\begin{proof}
Conditions~(1)--(5) collectively constitute the full support structure
of a branch book.  By Theorem~\ref{thm:dependency-declaration}, every
descendant theorem requires its support to be explicitly declared to
have canonical force.  The five items list exactly the support
categories relevant to a branch book.  The force restriction --- no
stronger status than declarations permit --- follows from
Proposition~\ref{prop:cite-declared-force}: a claim may only be cited
with the force its declared dependencies license.
\end{proof}

\begin{corollary}[\status{U}]\label{cor:no-alternate-foundation}
No branch book may act as an alternate foundation for NFC.
\end{corollary}

\begin{proof}
A branch book is a lawful descendant of $\mathbf{U}$ certified by
Book~V.  An alternate foundation would require providing a competing
universal source, which is impossible by Book~IV,
Theorem~IV.5.1 (uniqueness up to source equivalence).  Therefore any
branch book claiming foundational status is in violation of the spine
architecture.
\end{proof}

%% -- No-Hidden-Upgrade Rule ----------------------------------------

\begin{proposition}[\status{U}]\label{prop:no-hidden-upgrade}
\textup{(No-Hidden-Upgrade Rule.)}  A proposed endpoint label $L$ for
a branch book is admissible only if it passes all four tests:
\begin{enumerate}[label=\textup{(\roman*)}]
  \item\label{nhu:eliminable} \emph{Eliminability.}  $L$ is
        eliminable from every load-bearing statement in the branch
        book without affecting the mathematical content.
  \item\label{nhu:no-ontology} \emph{No ontology upgrade.}  $L$ does
        not assert an independent ontology beyond what the certified
        spine image supports.
  \item\label{nhu:no-enlargement} \emph{No source-image enlargement.}
        $L$ does not depend on structure outside the certified
        source image $F(\mathbf{U})$.
  \item\label{nhu:no-scope} \emph{No scope inflation.}  $L$ does not
        assert force at a scope wider than the branch book's declared
        domain of certification.
\end{enumerate}
A label failing any one of these tests is a prohibited import
(Definition~\ref{def:prohibited-import}) at the endpoint level.
\end{proposition}

\begin{proof}
Tests~\ref{nhu:eliminable}--\ref{nhu:no-scope} are the minimal
conditions under which an interpretive label adds no new mathematical
content to the branch and therefore inherits no more force than the
branch's theorems support.  Test~\ref{nhu:eliminable}: if $L$ cannot
be eliminated, it carries load-bearing mathematical content not
declared in the spine import ledger, violating
Theorem~\ref{thm:dependency-declaration}.
Test~\ref{nhu:no-ontology}: an independent ontology imports
structure outside the certified source image, violating the
no-smuggling discipline (Book~I, Standing~Rule~I.1.3).
Tests~\ref{nhu:no-enlargement} and~\ref{nhu:no-scope} are direct
consequences of the Canonical Closure Schema
(Theorem~\ref{thm:closure-schema}).
\end{proof}

\begin{remark}[\status{R}]
The No-Hidden-Upgrade Rule applies to every branch label, not only to
gravity labels.  ``Consciousness,'' ``spacetime,'' ``probability,''
``mass gap,'' and any other endpoint label must clear all four tests
before the branch book may employ it.  Failure is not a deficiency of
the mathematics; it is an honest naming violation that must be
corrected by relabeling, not by argument.
\end{remark}

\begin{remark}[\status{R}]
\textup{(v7.35 Smuggling Audit as governance template.)}
The pre-canonical program (v7.35 \S13206) maintains an explicit
Smuggling Audit table recording for each major result: whether new
primitives were introduced, whether the derivation is analytic or
computational, and whether the smuggling verdict is PASS or FAIL
(fields: $\{\mathrm{Result} \mid \mathrm{New\ primitive?} \mid
\mathrm{Derivation\ type} \mid \mathrm{Smuggling\ verdict} \mid
\mathrm{Justification}\}$).  This is a direct precursor to the
admissible-import/prohibited-import machinery of this book and the
No-Hidden-Upgrade Rule (Prop.~\ref{prop:no-hidden-upgrade}).  The
v7.35 table could serve as the template for systematic provenance
audits of future canonical additions.  Its format is compatible with
the Certificate Catalog Schema (Prop.~\ref{prop:cert-catalog}).
\end{remark}

%% -- Non-Retroactivity of Branch Certification --------------------

\begin{theorem}[\status{U}]\label{thm:non-retroactive}
\textup{(Non-Retroactivity of Branch Certification.)}  No newly
admitted branch may retroactively alter the burden, meaning, or force
of previously certified branches.
\end{theorem}

\begin{proof}
A previously certified branch book has been validated against the spine
through its declared import ledger, endpoint target, and named failure
frontier.  A new branch book is a new derivation from the same source
$\mathbf{U}$, certified independently through Book~V.  By
Theorem~\ref{thm:stable-extensibility}, the canonical spine is stable
under branch extension: new branches extend the program without
revising prior certified content.  Therefore the force, scope, and
burden of a previously certified branch are determined by its own
declarations, not by the content or closure of any later branch.
\end{proof}

\begin{remark}[\status{R}]
Non-retroactivity protects the honest-posture discipline: a branch
book that names an open frontier cannot have that frontier
retroactively closed by a later branch that claims a different
endpoint.  The NS Branch Book's Stage-3 frontier is its own named
obligation, independent of whether the YM Branch Book reaches
CERT-CLOSE.
\end{remark}

%% -- Shared Endgame Schema -----------------------------------------

\begin{proposition}[\status{U}]\label{prop:shared-endgame}
\textup{(Shared Endgame Schema.)}  Every branch book reaching its
claimed endpoint must establish the following three-part closure
structure:
\begin{enumerate}[label=\textup{(\roman*)}]
  \item \emph{Source-controlled realized object:}  A realized object
        $F(\mathbf{U})$ in the target category, descended from the
        universal source through a lawful realization functor.
  \item \emph{Branch-internal evolution law:}  A derivation, within
        the branch book, of the evolution law, gap law, or structural
        compatibility law appropriate to the endpoint.
  \item \emph{One-step interface propagation with branch-specific
        subcriticality:}  Certified one-step admissible enlargement
        that propagates the realized object's load-bearing structure,
        with the subcriticality condition verified branch-locally.
\end{enumerate}
For the Yang--Mills branch: parts~(i)--(iii) yield the mass gap.
For the Navier--Stokes branch: part~(ii) is the quasi-decay law and
part~(iii) is the Stage-3 interface bound.
For the gravity branch: part~(ii) is the deformation-response law and
part~(iii) is the curvature-subcritical extension.
\end{proposition}

\begin{proof}
The three-part structure is the canonical decomposition of what the
Canonical Closure Schema (Theorem~\ref{thm:closure-schema}) requires
of a branch book claiming endpoint force.  Part~(i) corresponds to
items~(1)--(2) of Theorem~\ref{thm:closure-schema} (spine imports and
branch inputs).  Part~(ii) corresponds to item~(3) (endpoint target
derivation).  Part~(iii) corresponds to items~(4)--(5) (bridge theorem
and named failure frontier if propagation fails).  Every branch
claiming an endpoint must provide all three; a branch lacking any part
has at most CERT-PROJ status.
\end{proof}

\begin{corollary}[\status{U}]\label{cor:undeclared-noncanonicalnical}
A theorem lacking declared import pathway has no canonical force.
\end{corollary}

\begin{proof}
By Theorem~\ref{thm:dependency-declaration}.
\end{proof}

%% ---------------------------------------------------------------
\section{Provisional Completion and External Readiness}
%% ---------------------------------------------------------------

The project must explicitly distinguish fixed architecture from
finished theorem content.

\begin{definition}[\status{D}]\label{def:architecturally-fixed}
A book is \emph{architecturally fixed} when its canonical role in the
seven-book spine is settled and will not be revised.
\end{definition}

\begin{definition}[\status{D}]\label{def:proof-hardening}
A book is \emph{under proof hardening} when its architectural role is
fixed but its theorem content is still being strengthened, split, or
rewritten within the fixed architecture.
\end{definition}

\begin{definition}[\status{D}]\label{def:externally-ready}
A book is \emph{externally ready} when its architecture is fixed, its
dependency ledger is explicit, its proofs have been hardened to the
current frontier, and its remaining open obligations are precisely
named.
\end{definition}

%% -- Status Correction Rule (I2) ----------------------------------

\begin{proposition}[\status{U}]\label{prop:status-correction}
\textup{(Status Correction Rule.)}  No canonical claim may have its
status upgraded --- from [C] to [U], from [O] to [C], or from any
weaker status to any stronger status --- without a formal ledger
revision event consisting of: (i)~a specific new proof or derivation
within the canonical spine or branch book; (ii)~an explicit declaration
of the upgraded claim's new status, scope, and dependencies in the
canonical paper; and (iii)~a formal acknowledgment that the prior
weaker status entry has been superseded.  Claims upgraded without a
ledger revision event retain their prior weaker status for canonical
purposes.
\end{proposition}

\begin{proof}
By Proposition~\ref{prop:cite-declared-force}, a claim may only be
cited with the force its declared dependencies license.  A status
upgrade without a formal canonical derivation provides no new canonical
proof; the prior declaration therefore continues to govern.
\end{proof}

%% -- Certificate Governance ---------------------------------------

\begin{proposition}[\status{U}]\label{prop:cert-catalog}
\textup{(Certificate Catalog Schema.)}  Every certificate deployed in
a canonical proof must be recorded in a certificate catalog with the
fields
\[
  \langle\mathrm{label},\, R,\, \Lambda,\, \Sigma,\,
  \delta,\, \mathrm{Inst}(\mathrm{label})\rangle
\]
(regime, predicate set, specification language, tolerance,
instantiation record).  A certificate without a catalog entry has
no canonical force.
\end{proposition}

\begin{proof}
By Theorem~\ref{thm:dependency-declaration}, every dependency must be
explicitly declared.  A certificate is a dependency; without a catalog
entry it cannot be independently verified or replicated, and therefore
lacks canonical force.
\end{proof}

\begin{proposition}[\status{U}]\label{prop:repository-law}
\textup{(Repository Law.)}  In the canonical text, certificates are
cited functionally by label.  Computation logs, trace files, and raw
output data belong in a repository companion, not in the theorem path.
\end{proposition}

\begin{proof}
The theorem path must be independently reformulable and extendable.
Raw computational output in the theorem path conflates the mathematical
claim with its evidentiary support and prevents clean re-derivation.
The Repository Law separates these layers: the theorem path cites only
the certificate label and its predicate-completeness status; the
repository companion holds the evidence.
\end{proof}

\begin{proposition}[\status{U}]\label{prop:software-invariance}
\textup{(Software-Invariance of Instantiation.)}  Two certificates for
the same regime with the same predicate set have identical admissible
instantiations.  The choice of computation software is invisible to
the mathematics.
\end{proposition}

\begin{proof}
Admissible instantiations are determined by the predicate set
$\Lambda$ and regime $R$.  Two certificates agreeing on $(\Lambda, R)$
agree on every mathematical consequence.  The software that produced
the certificate is not part of $\Lambda$ and cannot affect what the
certificate proves.
\end{proof}

\begin{theorem}[\status{U}]\label{thm:provisional-completion}
\textup{(Provisional Completion Principle.)}  The canonical
architecture of NFC may be fixed before the project is
theorem-level complete.  Theorem-level completion remains provisional
until the source spine has been hardened and the principal derived
branches have been pushed to their furthest honest closure prior to
external scrutiny.
\end{theorem}

\begin{proof}
Architectural fixity is a structural property: the seven books are
assigned their roles and will not be reorganized.  Theorem-level
completeness is a content property: every load-bearing theorem has
been proved or has a named frontier.  These are independent properties.
An architecture can be fixed --- preventing drift and enabling coherent
development --- while theorem content remains under active hardening.
Therefore the two need not coincide.
\end{proof}

\begin{corollary}[\status{U}]\label{cor:stability-not-finality}
Structural stability does not imply theorem-level finality.
\end{corollary}

\begin{remark}[\status{R}]\label{rmk:provisional-principle-canonical}
Theorem~\ref{thm:provisional-completion} is not merely project
management advice.  It is a canonical constitutional law protecting the
program from two opposite mistakes: endlessly reopening architecture
when only proofs need hardening, and claiming the work is complete
merely because the structure is settled.  Both errors have damaged
large mathematical programs.  The Provisional Completion Principle
prevents both.
\end{remark}

%% ---------------------------------------------------------------
\section{Canonical Delta Visibility and Replayability}
%% ---------------------------------------------------------------

The canon must remain governable over time.

\begin{definition}[\status{D}]\label{def:delta-visibility}
\emph{Canonical delta visibility} is the requirement that every
authoritative change to the canon be reconstructible from the
certified transition record: any addition, removal, or modification
of a canonical claim must be recorded explicitly.
\end{definition}

\begin{definition}[\status{D}]\label{def:replayability}
\emph{Canonical replayability} is the requirement that an independent
steward be able to reconstruct the present canonical state from the
declared transition record, beginning from the primitive definitions
of Book~I.
\end{definition}

\begin{proposition}[\status{U}]\label{prop:evolution-replayable}
Canonical evolution must remain replayable from the certified record
of changes.
\end{proposition}

\begin{proof}
By Definitions~\ref{def:delta-visibility} and \ref{def:replayability},
replayability requires that every canonical change be recorded.  If
changes occur without record, then an independent steward cannot
reconstruct the present state, and the canon has evolved in a way
invisible to the community.  The proposition is therefore a
consequence of the two definitions and the requirement that the canon
be externally auditable.
\end{proof}

\begin{corollary}[\status{U}]\label{cor:hidden-accretion-noncanonicalnical}
Authoritative canonical change without replayable record is
non-canonical.
\end{corollary}

\begin{remark}[\status{R}]\label{rmk:replayability-governance}
Proposition~\ref{prop:evolution-replayable} governs canon stewardship
rather than source mathematics.  It may ultimately be compressed into
an appendix note, but it belongs in Book~VII at least in concise form:
the constitutional book should record not only the laws governing
theorem force, but also the law governing how the constitution itself
may be amended.
\end{remark}

%% ---------------------------------------------------------------
\section{Governing Theorem of Book VII}
%% ---------------------------------------------------------------

\begin{theorem}[\status{U}]\label{thm:governing}
\textup{(Canonical Closure Schema Theorem of the Spine.)}  The
seven-book spine determines the exact rules under which descendant
branches may inherit, extend, and conditionally specialize certified
source structure, while preserving:
\begin{enumerate}[label=\textup{(\arabic*)}]
  \item explicit limits of force (status and dependency discipline),
  \item named open frontiers (failure-frontier stability),
  \item prohibition of hidden supplementation (import discipline),
  \item and prohibition of silent status inflation (promotion law).
\end{enumerate}
\end{theorem}

\begin{proof}
Combine:
\begin{itemize}
  \item the status and dependency discipline of Section~\ref{sec:status}:
        Proposition~\ref{prop:cite-declared-force},
        Corollary~\ref{cor:no-unconditional-by-proximity};
  \item the scoped-certificate and promotion law of
        Section~\ref{sec:promotion}:
        Theorems~\ref{thm:scoped-cert-rule},
        \ref{thm:promotion-requires-transfer};
  \item the failure-frontier stability of Section~\ref{sec:frontiers}:
        Theorem~\ref{thm:frontier-stability};
  \item the import discipline and closure schema of Section~VI:
        Theorems~\ref{thm:dependency-declaration},
        \ref{thm:closure-schema};
  \item and the extensibility and falsifiability conditions of
        Section~\ref{sec:extensibility}:
        Theorems~\ref{thm:stable-extensibility},
        \ref{thm:external-falsifiability}.
\end{itemize}
Together these establish the full constitutional framework governing
how branch books may use, cite, extend, and close results from the
canonical spine.
\end{proof}

\begin{corollary}[\status{U}]\label{cor:noncompliant-noncanonicalnical}
Any descendant claim lacking explicit status, imports, scope, and
frontier discipline is non-canonical.
\end{corollary}

\begin{scholium}[\status{R}]\label{sch:VII-indispensable}
Book~VII adds no new source apex, no new branch endpoint, and no new
continuum interface.  It determines how all such content may be
lawfully claimed.  That is why the book is indispensable.  A
mathematical program without an explicit theory of its own epistemic
standards is not a rigorous program; it is a collection of results
waiting to be challenged.  Book~VII provides those standards as
canonical law.
\end{scholium}

%% ---------------------------------------------------------------
%% ---------------------------------------------------------------
\section{Canonical Program Status Proposition}
\label{sec:program-status}
%% ---------------------------------------------------------------

\begin{proposition}[\status{U}]\label{prop:program-status}
\textup{(NFC Canonical Program Status.)}
The Nested Fibrational Cosmology canonical program currently satisfies:
\begin{enumerate}[label=\textup{(\Roman*)}]
  \item \textbf{Seven-book spine}: Books~I--VII compile without
    LaTeX errors; all eleven canonical files compile cleanly.
    The spine is internally consistent under the declared
    governance rules.

  \item \textbf{Conditionally closed branches}:
    \begin{itemize}
      \item \emph{YM Branch}: Clay-grade conditional CERT-CLOSE
        achieved --- mass gap $m^2 > 0$ globally on $\mathbb{R}^4$
        with OS/Wightman axioms, conditional on Phase-A, $K_0=7$,
        gap-subcriticality, and KPO-3.  Three post-program gaps
        B1/B2/B3 explicitly registered.
      \item \emph{NS Branch}: domain-bounded CERT-CLOSE --- Leray
        weak solution $u \in L^\infty_t H^1_x$ conditional on
        NS.6.1 hypotheses and window-uniformity.
      \item \emph{SCC Branch}: conditional CERT-CLOSE chain
        (WeakGlue $\Rightarrow$ MCS $\Rightarrow$ UC $\Rightarrow$
        TSI) fully discharged conditionally.
      \item \emph{GR Branch}: extended domain-bounded CERT-CLOSE
        under curvature-subcriticality; EFE continuum limit
        theorem conditional on KPO-3.
    \end{itemize}

  \item \textbf{Key structural resolutions}:
    \begin{itemize}
      \item E$_8$/O-ID conflict conditionally resolved via
        $d_{S_v}$-block screening: both routes reconciled at
        the level of structural vs.\ observable algebra.
      \item PASS is a theorem (not an axiom): derived from
        Phase-A + pendant gauge covariance.
      \item $K_0 = 7$ is the unique prime fixed point of
        $F(p) = |\hat\Phi(\mathfrak{g}_{\mathrm{color}}(p))|$:
        derived unconditionally.
      \item Lorentzian signature $(-,+,+,+)$ emergent from
        Phase-A + spatial isotropy (conditional).
      \item $\beta = \log_2(3/2)$: canonical defect-ledger unit
        (unconditional).
    \end{itemize}

  \item \textbf{RH program}: 9 of 10 admissible-capture
    obstructions conditionally discharged.  One remains:
    S1 (witness assembly) --- requires arithmetic witness
    construction from the number-theoretic zero descent of $\Xi$.

  \item \textbf{Open obligations} (honest frontier):
    \begin{itemize}
      \item RH S1 (witness assembly) --- arithmetic/number theory;
      \item GR global curvature-subcriticality --- nonlinear
        stability of the maximal EFE development;
      \item YM B1/B2/B3 (quantization, domain, gauge group
        simplicity) --- post-program, out of NFC scope;
      \item O-ID\textsuperscript{cont} $O(n^{-1})$ rate
        certification;
      \item NS window-uniformity verification (the structural
        hypothesis in NS.6.1 global promotion).
    \end{itemize}

  \item \textbf{Governance}: all promotions satisfy the three
    readiness criteria (precise canonical address, no unlicensed
    new primitives, proof or declared conditional available);
    all conflicts are named and either resolved or quarantined;
    all open obligations carry their precise discharge conditions.
\end{enumerate}
\end{proposition}

\begin{proof}
Each claim is a summary of the established canonical content
and is independently verifiable by inspecting the cited
theorems and their status tags.  Book~VII's governance
structure guarantees that no silent promotion has occurred:
every [C] item carries its conditions, every [O] item names
its discharge criterion, and every [U] item is proved. \qed
\end{proof}

%% ---------------------------------------------------------------
\section{Branch Status Table}
\label{sec:status-table}
%% ---------------------------------------------------------------

\begin{remark}[\status{R}]
\textbf{Complete branch status as of the current canonical state.}

\renewcommand{\arraystretch}{1.3}
\begin{center}
\begin{tabular}{|l|l|l|}
\hline
\textbf{Branch / Book} & \textbf{Status} & \textbf{Key Remaining} \\
\hline
Book I (Primitive Foundations) & [U] Stable & --- \\
Book II (Local Structure) & [U]/[C] Stable & --- \\
Book III (Transport/Persistence) & [U]/[C] Stable & --- \\
Book IV (Universal Source) & [U] Stable & --- \\
Book V (Descent/Legitimacy) & [U]/[C] Stable & --- \\
Book VI (Continuum Interface) & [U]/[C] Stable & --- \\
Book VII (Governance) & [U]/[D] Stable & --- \\
\hline
YM Branch & Cond.\ CERT-CLOSE & B1/B2/B3 post-program \\
NS Branch & Cond.\ CERT-CLOSE & LR-weighted global \\
SCC Branch & Cond.\ CERT-CLOSE & UCTI certification \\
GR Branch & Cond.\ CERT-CLOSE & Global curv.-subcrit. \\
\hline
SM Branch & CERT-PROJ & O-SM.Matter, CouplingTransfer \\
RH Branch & CERT-PROJ & RH4--6, S1 (arithmetic) \\
\hline
\end{tabular}
\end{center}

\textbf{Irreducible open frontier} (outside the canonical toolkit):
\begin{itemize}
  \item \textbf{RH S1}: Arithmetic witness assembly (Mellin/Dirichlet
    descent from $\Xi$) --- pure number theory.
  \item \textbf{GR global curvature-subcriticality}: Nonlinear
    stability of the maximal EFE Cauchy development.
  \item \textbf{SM O-SM.Matter}: Quarks, leptons, Higgs field content.
  \item \textbf{SM O-SM.CouplingTransfer}: Coupling constant derivation
    from archetype interface sharing (highest-value open obligation).
\end{itemize}

\textbf{Post-program} (explicitly out of NFC scope):
\begin{itemize}
  \item YM B1: Classical vs.\ quantum Hamiltonian.
  \item YM B2: Abstract algebra domain vs.\ $\mathbb{R}^4$.
  \item YM B3: Reductive SM gauge group vs.\ compact simple $G$.
\end{itemize}
\end{remark}



%% ---------------------------------------------------------------
\section{Named Obligations vs.\ Reduced Irreducible Frontier}
\label{sec:obligation-vs-frontier}
%% ---------------------------------------------------------------

\begin{remark}[\status{R}]\label{rem:named-vs-frontier}
\textbf{Two distinct ledger layers.}
The canonical program maintains two distinct notions that must not be
conflated:

\begin{enumerate}[label=\textup{(\arabic*)}]
  \item \textbf{Named open obligations} $[\mathrm{O}]$ --- formal
    declarations inside canonical files with named labels, stating
    exactly what remains unproved at the current canonical stage.
    These are the governance document.  Every named obligation in the
    canonical files (currently 13 items) has a label, a declared
    scope, and an auditable discharge criterion.

  \item \textbf{Reduced irreducible frontier} --- the smallest
    set of genuinely open mathematical problems remaining after all
    hardening passes, structural reductions, conditional discharges,
    and iff-schema decompositions have been applied.  The frontier
    is always a subset of the named obligations, but it may be
    strictly smaller.
\end{enumerate}

\textbf{Why they differ.}
A named obligation may be:
\begin{itemize}
  \item \emph{Post-program} (explicitly out of NFC scope): named
    to record the Clay distance honestly, not as an active discharge
    target.  Example: ob:YM-B1/B2/B3 (classical vs.\ quantum
    Hamiltonian; domain; gauge group).
  \item \emph{Gated} (becomes tractable only once an upstream
    obstacle is resolved): named to make the dependency visible.
    Example: ob:O-ID-cont (gated on the E$_8$/SM arena separation);
    ob:SM-harmonization and ob:SM-coupling-transfer variants (gated
    on O-SM.Matter content).
  \item \emph{Hardened} (the obligation label persists but its
    content has been sharply reduced by subsequent work): the label
    did not disappear, but what it means is now much more precise.
    Example: ob:rh-d1-full, which after the full RH witness ladder
    (Thms.~\ref{thm:rh-f3}--\ref{thm:rh-f4}, RH Branch) is now
    equivalent to seven named uniform laws, with the two deepest
    hinges identified explicitly.
  \item \emph{Duplicated in formulation}: two obligation labels
    may target the same mathematical gap at different levels of
    precision.  Example: ob:SM-coupling-transfer and
    ob:SM-coupling-transfer-full.
\end{itemize}

\textbf{The reduced irreducible frontier} (post-hardening, as of the
current canonical state) consists of exactly four items:
\begin{enumerate}[label=\textup{(\arabic*)}]
  \item \emph{Arithmetic orbit grammar on $\Xi$}: the first hard
    block.  The scaling-flow candidate $G_{\mathrm{sf}}$ is now
    conditionally in place: the orbit grammar
    (Def.~\ref{def:rh-orbit-grammar}, RH Branch), the D2 local
    audit (Thm.~\ref{thm:rh-sf-local-d2-checks}), and the RH-N2
    realization (Cor.~\ref{cor:rh-sf-first-hard-block-cleared})
    are all conditionally discharged.  The first hard block is
    cleared for the scaling-flow candidate pending full
    branch-internal proof of Defs.~\ref{def:rh-logdiff-channel}
    and~\ref{def:rh-trace-test-subfamily}.
  \item \emph{Packet-local synthesis law}: the second hard block.
    Steps L1 (packet decomposition) and L2 (packet-scale
    extraction) are conditionally discharged
    (Props.~\ref{prop:rh-tloc-L1}--\ref{prop:rh-tloc-L2}, RH
    Branch).  The live core is L3 (phase detectability), L4
    (phase-spacing), L5--L7 (Gram, leakage, compression).
    When all seven steps are discharged,
    Thm.~\ref{thm:rh-wc1bb} follows.
  \item \emph{SM matter content}: conditionally advanced at
    structural level.  Source-descended fermionic representations
    (Thm.~\ref{thm:SM-matter-source-descended}, SM Branch) and
    $B_0$ Higgs mechanism (Thm.~\ref{thm:SM-higgs-source-descended})
    are established.  Full coupling ratio $g_3^2:g_2^2:g_1^2 = 6:2:1$
    proved (Thm.~\ref{thm:SM-coupling-transfer-full}).
    Remaining gap: O\_ID\textsuperscript{cont} continuum
    identification at full canonical force for the complete
    physical-sector construction.
  \item \emph{GR global curvature-subcriticality}: nonlinear EFE
    stability beyond the domain-bounded regime --- the extension of
    the domain-bounded CERT-CLOSE to a globally valid statement
    (def:curv-subcrit-global, GR Branch).
\end{enumerate}

Everything else on the 13-item named-obligation list is either
post-program (YM B1/B2/B3), gated (ob:O-ID-cont, ob:SM-harmonization,
coupling formulations), or hardened-but-relabeled versions of the
four items above.

\textbf{Governance rule.}
This distinction is itself a canonical record.  A named obligation
may not be claimed as ``discharged'' merely because its frontier
content has been reduced to a sharper formulation --- discharge
requires a theorem with proof at the declared level.  Conversely,
the hardening of a named obligation into a precise iff-schema
is positive progress and must be recorded in the Canon Ledger,
even when the obligation label is retained.

\textbf{Reading the 13-item snapshot.}
The 13 canonical `[O]' items are the correct governance document for
auditing the program.  The four-item reduced frontier above is the
correct document for assessing the genuine mathematical distance to
completion.  Both are needed; neither replaces the other.
\end{remark}

\begin{remark}[\status{R}]\label{rem:obligation-conventions}
\textbf{Counting conventions and the authoritative live record
(governance hardening).}
Two syntactic conventions for recording obligations coexist in the
canon, and they must be read together, not mechanically summed:
\begin{enumerate}[label=\textup{(\roman*)}]
  \item \emph{Status-tagged remarks.}  An obligation written as a
    remark or declaration carrying the $[\mathrm{O}]$ status tag.
    This is the Book~I--III and branch-frontier convention; a
    syntactic scan for $[\mathrm{O}]$ finds exactly these.
  \item \emph{The \texttt{openobligation} environment.}  The YM
    branch (and any branch adopting its convention) records
    obligations in a dedicated \texttt{openobligation} environment
    that is \emph{status-tagged} $[\mathrm{C}]$ when the obligation
    is conditionally discharged at a named scope --- e.g.\ the YM
    Clay gaps ob:YM-B1/B2/B3, conditionally discharged at NFC
    scope with the residual to full quantum field theory named.
    A syntactic scan for $[\mathrm{O}]$ does \emph{not} find these,
    because they are honestly $[\mathrm{C}]$, not open-unproved.
\end{enumerate}
\textbf{Consequence for census hygiene.}  A raw $[\mathrm{O}]$-tag
count (the syntactic census) and the curated \emph{named-obligation}
count of this section (the semantic governance count) measure
different things and need not coincide; likewise the curated
\emph{reduced irreducible frontier} is a mathematical-distance
assessment on a third basis again.  None of the three is wrong, but
none is a substitute for another, and a reader who reports a single
``frontier size'' without naming the basis will mislead.
\textbf{Snapshot lag is expected.}  The curated counts in this
section are periodic snapshots requiring manual update; they may lag
the syntactic state after an active arc.  The current known lag: the
Variable Recovery Program arc (Book~III, recorded in the Canon
Ledger) opened and then discharged or characterized its full
dynamical-core obligation set, which this section's ``13-item''
snapshot predates and does not yet reflect.  \textbf{Authoritative
record.}  The Canon Ledger is the authoritative live record of
obligation state, session by session; the snapshots here are
governance summaries reconciled to it only at dedicated
reconciliation passes.  A reconciliation pass --- recomputing the
named-obligation roster and the reduced frontier against the current
syntactic state across all books --- is itself a registered
governance task, distinct from the per-session ledger discipline,
and is not performed implicitly by branch work.  This remark records
the convention and the standing reconciliation task; it changes no
count, since responsible recounting requires the dedicated pass.
\emph{Status update:} that reconciliation pass was carried out and is
recorded in the external governance document
\texttt{NFC\_OBLIGATION\_ROSTER} --- the authoritative partition of all
109 corpus \texttt{ob:} labels (open / conditionally-discharged /
bridge-discharged / fully-discharged / base-declaration), with the
reduced four-item frontier re-verified against current live labels and
four tracking conventions catalogued.  The roster supersedes the
``13-item'' figure of Rmk.~\ref{rem:named-vs-frontier} as the current
snapshot, while that remark's \emph{distinction} between named
obligations and the reduced irreducible frontier stands unchanged: the
four-item reduced frontier was re-verified and remains correct, with
the GR item noted to be tracked as \texttt{def:curv-subcrit-global}
rather than an \texttt{ob:} label.  The Canon Ledger remains the live
record between snapshots.
\end{remark}

\section{Boundary of Book VII}
%% ---------------------------------------------------------------

\textbf{What Book~VII proves.}
\begin{enumerate}
  \item Every theorem must declare its status, dependencies, and
        scoped regime.
  \item No conditional result becomes unconditional by repetition or
        proximity.
  \item Scoped certificates carry force only in their declared regime.
  \item Promotion requires both a scoped certificate and a separately
        proved transfer theorem.
  \item Status inflation --- claiming more than declared dependencies
        license --- is prohibited.
  \item Named failure frontiers are positive constitutional objects;
        rephrasing does not bypass them.
  \item The canonical spine is stably extensible and every lawful
        branch is externally falsifiable.
  \item Every branch book must declare its five support components
        and may claim no stronger status than those declarations permit.
  \item No branch book may act as an alternate foundation for NFC.
  \item Structural stability does not imply theorem-level finality.
  \item \emph{New enrichments.}  The No-Hidden-Upgrade Rule
        (Prop.~\ref{prop:no-hidden-upgrade}): four-test admissibility
        for any endpoint label.  Non-Retroactivity of Branch
        Certification (Thm.~\ref{thm:non-retroactive}): no new branch
        alters prior certified force.  The Shared Endgame Schema
        (Prop.~\ref{prop:shared-endgame}): the three-part closure
        structure required for any branch claiming an endpoint.  The
        Status Correction Rule (Prop.~\ref{prop:status-correction}).
        Certificate governance: catalog schema, repository law, and
        software-invariance (Props.~\ref{prop:cert-catalog}--\ref{prop:software-invariance}).
\end{enumerate}

\textbf{What Book~VII does not govern.}
\begin{enumerate}
  \item \emph{Mathematical content of specific branches.}  Book~VII
        governs how branch claims must be stated and cited; it does not
        determine whether any specific branch claim is true.
  \item \emph{Resolution of named open obligations.}  Book~VII
        identifies the frontier structure; it does not close it.  That
        remains the work of branch books and future theorem-hardening.
  \item \emph{Priority or ordering among branches.}  Book~VII treats
        all lawful branches equally under the same governance law.
\end{enumerate}

%% ---------------------------------------------------------------
\section{Dependency Ledger}
%% ---------------------------------------------------------------

\renewcommand{\arraystretch}{1.4}
\noindent
\small
\begin{tabular}{@{}p{4.4cm}p{0.8cm}p{3.0cm}p{4.8cm}@{}}
\hline
\textbf{Theorem} & \textbf{St.} & \textbf{Depends on} & \textbf{Exports to} \\
\hline
\ref{prop:cite-declared-force}\ Cite declared force
  & [U] & SR~\ref{sr:declare-status}, Def.~\ref{def:unconditional}
  & \ref{cor:no-unconditional-by-proximity}, \ref{thm:governing} \\

\ref{thm:scoped-cert-rule}\ Scoped Certificate Rule
  & [U] & Def.~\ref{def:scoped-cert}
  & \ref{thm:promotion-requires-transfer}, \ref{cor:finite-no-continuum} \\

\ref{thm:promotion-requires-transfer}\ Promotion Req.\ Transfer
  & [U] & \ref{thm:scoped-cert-rule}, Def.~\ref{def:transfer-theorem}
  & \ref{cor:finite-no-continuum}--\ref{cor:conditional-no-unconditional};
    all branch books \\

\ref{prop:downstream-force}\ Downstream force
  & [U] & \ref{prop:cite-declared-force}
  & \ref{cor:failed-bridge-reopens}, \ref{cor:slogan-no-force} \\

\ref{thm:frontier-stability}\ Named Failure Frontier Thm.
  & [U] & Def.~\ref{def:named-frontier}
  & \ref{cor:rephrasing-not-progress}; all branch books \\

\ref{thm:stable-extensibility}\ Stable Extensibility Thm.
  & [U] & Bk.~V branch law, Defs.~\ref{def:stably-extensible}
  & \ref{thm:governing}, governance \\

\ref{thm:external-falsifiability}\ External Falsifiability Thm.
  & [U] & Def.~\ref{def:externally-falsifiable}
  & \ref{cor:opaque-inadmissible}; all branch books \\

\ref{thm:dependency-declaration}\ Dependency Declaration Thm.
  & [U] & \ref{prop:cite-declared-force},
    Def.~\ref{def:admissible-import}
  & \ref{thm:closure-schema}, \ref{cor:undeclared-noncanonicalnical} \\

\ref{thm:closure-schema}\ Canonical Closure Schema Thm.
  & [U] & \ref{thm:dependency-declaration},
    \ref{prop:cite-declared-force}
  & \ref{cor:no-alternate-foundation}, \ref{thm:governing};
    all branch books \\

\ref{thm:provisional-completion}\ Provisional Completion Princ.
  & [U] & Defs.~\ref{def:architecturally-fixed}--\ref{def:externally-ready}
  & \ref{cor:stability-not-finality}; project governance \\

\ref{prop:no-hidden-upgrade}\ No-Hidden-Upgrade Rule
  & [U] & \ref{thm:dependency-declaration};
    Bk.~I SR~I.1.3 (no-smuggling); \ref{thm:closure-schema}
  & All branch books (endpoint label discipline) \\

\ref{thm:non-retroactive}\ Non-Retroactivity
  & [U] & \ref{thm:stable-extensibility}; Bk.~V cert.
  & All branch books; multi-branch programs \\

\ref{prop:shared-endgame}\ Shared Endgame Schema
  & [U] & \ref{thm:closure-schema}; Bks.~III--V
  & YM + NS + GR branch books Part~IV \\

\ref{prop:status-correction}\ Status Correction Rule
  & [U] & \ref{prop:cite-declared-force}
  & All branch books; ledger governance \\

\ref{prop:cert-catalog}\ Certificate Catalog Schema
  & [U] & \ref{thm:dependency-declaration}
  & All branch books; evidentiary layer \\

\ref{prop:repository-law}\ Repository Law
  & [U] & \ref{prop:cert-catalog}
  & All branch books; project management \\

\ref{prop:software-invariance}\ Software-Invariance
  & [U] & \ref{prop:cert-catalog}
  & All branch books; computation citations \\

\ref{prop:evolution-replayable}\ Evolution replayable
  & [U] & Defs.~\ref{def:delta-visibility}, \ref{def:replayability}
  & \ref{cor:hidden-accretion-noncanonicalnical}; governance \\

\ref{thm:governing}\ Canonical Closure Schema of Spine
  & [U] & All of Bk.~VII
  & Constitutional law for all branch books; \ref{cor:noncompliant-noncanonicalnical} \\
\hline
\end{tabular}
\normalsize
\renewcommand{\arraystretch}{1}

\medskip

\begin{scholium}[\status{R}]\label{sch:VII-final}
Book~VII closes the canonical spine.  Together, Books~I--VII
establish:
\begin{itemize}
  \item the primitive observational grammar (Book~I),
  \item local boundary control (Book~II),
  \item persistence and transport (Book~III),
  \item the universal source $\mathbf{U}$ (Book~IV),
  \item lawful descent to branches (Book~V),
  \item interface legitimacy for continuum tools (Book~VI), and
  \item constitutional force-governance for all of the above (Book~VII).
\end{itemize}
No further spine book is needed.  Derived branch books --- for
Yang--Mills, Navier--Stokes, gravity, and other endpoints --- may now
proceed under the architecture thus established, each certifying its
descent from $\mathbf{U}$ and declaring its claims under the
governance law of this book.
\end{scholium}


%% ---------------------------------------------------------------
\section{Terminal Status: Source Architecture and Residual Frontier}
\label{sec:terminal-status}
%% ---------------------------------------------------------------

\begin{remark}[\status{R}]\label{rem:terminal-source-reading}
\textbf{Source-centered reading of the corpus.}
The NFC program now carries a complete source-centered architecture
(dream\_paper\_v38.pdf, terminal edition):
\begin{itemize}
  \item \emph{Layer I (Unconditional):} the finite relational kernel
    establishes behavioral congruence, singular projection interfaces,
    finite stabilization, and defect bookkeeping from axioms A1--A4
    and canonical generator class L1--L4 alone.
  \item \emph{Layer II ($K_0 = 7$ as theorem):} $K_0$ must be prime
    (Thm.~\ref{thm:K0-prime}, Book~I [U]); $K_0 = 7$ is the unique
    prime fixed point of the Weyl structure identity
    $F(p) = |\hat\Phi(\mathfrak{g}_{\mathrm{color}}(p))|$ satisfying
    all three structural self-consistency conditions (F1)--(F3).
    This is not an assumption; it is a theorem of axioms A1--A4.
  \item \emph{Layer III (Conditional on $K_0 = 7$):} the universal
    object $\mathbf{U} \in \mathrm{POT}$
    (Thm.~\ref{thm:universal-completion}, Book~IV [U]); the unified
    coercive-transport source theorem BR
    (Thm.~\ref{thm:BR-source}, Book~V [C]); the ICDC compression
    schema (Thm.~\ref{thm:ICDC}, Book~V [C]); the gauge
    identification $G = \mathrm{SU}(3) \times \mathrm{SU}(2) \times
    \mathrm{U}(1)$; the controlled continuum realization; and the
    three branch specializations (YM, NS, GR).
\end{itemize}
The key structural fact is that the YM, NS, and GR branch theorems
are \emph{certified specializations of one upstream source theorem}
(BR), not independent foundations.  The terminal dependency order is:
finite source structure $\to$ universal source theorem $\to$ licensed
realization $\to$ branch-specific theorems.
\end{remark}

\begin{theorem}[\status{U}]\label{thm:anti-smuggling-boundary}
\textup{(Anti-Smuggling Boundary.)}
Let $C$ be a soundly scoped certificate for a claim $P$ on a
declared regime $R$.  Then $C$ establishes $P$ on $R$ only.  It
does not establish $P$ on any $R' \neq R$, nor any asymptotic,
continuum, or physical-interpretive extension of $P$ beyond $R$.
Three illegitimate moves are blocked: (i) treating a finite
computation as an asymptotic result; (ii) treating a certified
bounded regime as covering an unbounded one; (iii) silently upgrading
a finite check into a general theorem by rhetorical compression.
(Attributed: dream\_paper\_v38.pdf Thm.~5.2.)
\end{theorem}

\begin{proof}
A scoped certificate proves only the claim on its declared regime.
The Certificate Doctrine (Def.~\ref{def:scoped-certificate},
Book~VII) and the Continuum Bridge Schema
(Thm.~\ref{thm:continuum-bridge-schema}, Book~VI) jointly block
each of the three illegitimate moves: (i) follows from the
no-asymptotic-promotion standing rule; (ii) follows from the
anti-regime-expansion rule; (iii) follows from the no-rhetorical-upgrade
rule.  Hence $C$ has exactly the force declared by its scope. \qed
\end{proof}

\begin{proposition}[\status{U}]\label{prop:what-corpus-does-not-prove}
\textup{(Residual Frontier: What the Corpus Does Not Prove.)}
The NFC canonical corpus, as of the current revision, does not prove:
\begin{enumerate}[label=\textup{(\arabic*)}]
  \item \emph{Unconditional Millennium solutions.}  The YM, NS, and
    GR branch closures are conditional on $K_0 = 7$ and the
    licensed realization/bridge stack.  Rejecting any bridge
    reopens the corresponding downstream claim.
  \item \emph{Kernel-internal continuum results.}  The universal
    object, Haag--Kastler realization, and physical specializations
    depend on the Layer~III representation stack; they are not
    finite-kernel theorems in the narrow sense.
  \item \emph{Unconstrained gauge-sector uniqueness.}  The SM gauge
    identification $\mathrm{SU}(3) \times \mathrm{SU}(2) \times
    \mathrm{U}(1)$ holds inside the licensed bridge stack; it does
    not prove that every continuum interpretation of the kernel must
    yield the same group.
  \item \emph{Global curvature-subcriticality.}  The GR branch is
    domain-bounded (Thm.~\ref{thm:gr-domain-bounded}); global
    nonlinear EFE stability beyond $\mathcal{D}$ is an open
    obligation.
  \item \emph{Unconditional RH.}  The critical-line localization is
    conditional on ob:RH4 and ob:RH5 at full canonical force, and on
    extension beyond Phase-A scope.
  \item \emph{Clay-Prize-level YM.}  ob:YM-B1/B2/B3 record the
    honest Clay Math Prize distance; these are not intended for
    discharge within the current NFC program.
  \item \emph{Quantum gravity.}  No quantum-gravity theorem is
    claimed; it is a post-schema program admissible only after the
    gravity branch satisfies the shared endgame schema
    (Prop.~\ref{prop:gr-endgame-schema}, GR Branch).
\end{enumerate}
These are boundaries of force, not failure modes.  Every downstream
claim carries exactly the strength licensed by its named bridge stack.
\end{proposition}


%% ---------------------------------------------------------------
\section{CMP-Standard Audit}
\label{sec:cmp-audit}
%% ---------------------------------------------------------------

\begin{remark}[\status{R}]\label{rem:cmp-audit}
\textbf{CMP-Standardization audit (NFC\_-\_CMP-Standard.pdf, 805 pages).}
The CMP Standard document defines six program-wide requirements
and a 20-module proof-grade matrix (P0--P4, where P3 = CMP-like
internal standard, P4 = externally paper-ready).  The current corpus
satisfies all six requirements:
\begin{enumerate}[label=\textup{(\arabic*)}]
  \item \emph{Every object has a precise ambient home.}  Verified:
    all definitions include explicit ambient spaces (vector spaces,
    categories, algebras, Hilbert spaces as appropriate).
  \item \emph{Every theorem has a sharp status.}  Verified:
    the [U]/[C]/[D]/[O] tag system is enforced throughout;
    every conditional theorem states exact hypotheses, dependencies,
    scoped regime, and named frontier when not unconditional.
  \item \emph{Every selection principle is formalized.}  Verified:
    transport invariance, persistence-simple criterion, balanced
    residual exclusion, coercive-transport inequality all carry
    explicit definitions and proved lemmas.
  \item \emph{Branch-neutral vs.\ branch-specific separation.}
    Verified: the source/branch-neutral/branch-specific/continuum/endpoint
    architecture is enforced by the canonical label discipline and
    Book~VII governance.
  \item \emph{Heuristic ladders split into proved core / conjectural
    / roadmap.}  Verified: all speculative content is in the
    Speculative Holding file; all [C] theorems carry explicit
    conditional dependencies.
  \item \emph{Every major result has a referee attack surface.}
    Verified: Prop.~\ref{prop:what-corpus-does-not-prove} names
    the seven boundaries of force, and every branch closure remark
    states what a referee who rejects a bridge gets.
\end{enumerate}
\end{remark}

\begin{proposition}[\status{U}]\label{prop:cmp-p-levels}
\textup{(CMP Proof-Grade Assessment.)}
The 20-module CMP proof-grade matrix (P0 = vision only,
P3 = CMP-like internal standard) yields the following
honest P-level assignments for the current corpus:

\medskip
\noindent
\begin{tabular}{ll}
\hline
\textbf{Module} & \textbf{P-level} \\
\hline
Books I--III core (axioms, quotients, transport, UCTI) & P3.2 [U] \\
$K_0 = 7$ as theorem (primality + uniqueness) & P3.5 [U] \\
Persistence-simple + balanced residual exclusion & P3.0 [C] \\
E$_8$ graph reduction + Cartan realization & P3.2 [C] \\
Chevalley--Serre reconstruction + certificate & P3.4 [C] \\
Y2 derived continuum eligibility & P2.8 [C] \\
Y3/Y4/Y4.5 gauge precursor + transformation & P2.8 [C] \\
Y5/Y5.5 variational emergence + action uniqueness & P2.7 [C] \\
Y6 coercive mass-gap chain & P2.5 [C] \\
Universal object + BR source theorem & P3.0 [U/C] \\
ICDC schema + branch specializations & P3.0 [C] \\
SM gauge + 3 generations + matter content & P3.0 [C] \\
SM coupling $6:2:1$ Transfer Theorem & P2.8 [C] \\
Higgs $B_0$ screening mechanism & P2.5 [C] \\
RH orbit grammar + first hard block & P2.7 [C] \\
RH second hard block (synthesis template L1--L7) & P2.5 [C] \\
NS branch (Ren.5 + SB2 + conditional near-closure) & P3.0 [C] \\
GR branch (EFE structural + domain-bounded) & P2.5 [C] \\
Book VII governance & P3.5 [U] \\
\hline
\end{tabular}

\medskip
The largest upgrade from the CMP document's original projections:
Y6 mass gap (P0.5 $\to$ P2.5, the coercive five-lemma chain with
explicit $\varepsilon_0 > 0$), Y5 variational emergence
(P1 $\to$ P2.7), Y3/Y4 (P1.5 $\to$ P2.8), and Y2 continuum
eligibility (P1.5 $\to$ P2.8, with H4 derived from
transport-invariance + locality + persistence-simple minimality,
not assumed).
\end{proposition}

\begin{remark}[\status{R}]\label{rem:cmp-residual-gaps}
\textbf{Remaining CMP-grade gaps.}
The following items remain below P3 and represent the
honest residual frontier before full CMP-grade certification:
\begin{enumerate}[label=\textup{(\arabic*)}]
  \item \emph{Global nonlinear EFE stability} (GR):
    domain-bounded result is P2.5; global extension requires new
    analytic content outside the current NFC collar-algebra toolkit.
  \item \emph{RH beyond Phase-A scope}: the first/second hard block
    programs are at P2.5--P2.7 for the scaling-flow candidate; full
    unconditional RH requires lifting the Phase-A restriction.
  \item \emph{$O\_ID^{\mathrm{cont}}$ identification}: the full
    continuum identification of the screened sector (ob:O-ID-cont)
    would lift Y3--Y6 from P2.x to P3.x.
  \item \emph{Clay-Prize-level Yang--Mills} (ob:YM-B1/B2/B3):
    unconditional, beyond the current NFC program.
\end{enumerate}
These four items are the honest CMP-grade distance from P4
(externally paper-ready) for the most technically demanding modules.
\end{remark}

\end{document}
